% ---------------------------------------------------------------------------
% Informations à modifier avant remise du rapport
% ---------------------------------------------------------------------------
\newcommand{\Auteur}{Prénom NOM}
\newcommand{\Formation}{Intitulé de la formation}
\newcommand{\Etablissement}{Nom de l'établissement}
\newcommand{\Entreprise}{SNCF}
\newcommand{\EncadrantEntreprise}{Prénom NOM}
\newcommand{\EncadrantAcademique}{Prénom NOM}
\newcommand{\PeriodeStage}{Dates du stage}
\newcommand{\AnneeUniversitaire}{2025--2026}
\newcommand{\TitreRapport}{Conception d'un assistant GraphRAG pour le métamodèle Arianes}

\usepackage[utf8]{inputenc}
\usepackage[T1]{fontenc}
\usepackage[french]{babel}
\usepackage{lmodern}
\usepackage{microtype}
\usepackage{geometry}
\usepackage{graphicx}
\usepackage{float}
\usepackage{booktabs}
\usepackage{tabularx}
\usepackage{longtable}
\usepackage{array}
\usepackage{enumitem}
\usepackage{xcolor}
\usepackage{amsmath}
\usepackage{amssymb}
\usepackage{listings}
\usepackage{tikz}
\usetikzlibrary{arrows.meta,positioning,fit,shapes.geometric}
\usepackage{caption}
\usepackage{subcaption}
\usepackage{fancyhdr}
\usepackage{hyperref}

\geometry{top=2.5cm,bottom=2.5cm,left=2.7cm,right=2.3cm}
\setlength{\parindent}{0pt}
\setlength{\parskip}{0.65em}
\setlist[itemize]{topsep=0.25em,itemsep=0.2em}
\setlist[enumerate]{topsep=0.25em,itemsep=0.2em}
\renewcommand{\arraystretch}{1.2}

\definecolor{sncfviolet}{RGB}{100,45,110}
\definecolor{sncfrouge}{RGB}{210,0,55}
\definecolor{sncfgris}{RGB}{75,78,82}
\definecolor{fondcode}{RGB}{247,247,247}
\definecolor{vertvalidation}{RGB}{30,120,75}

\hypersetup{
  colorlinks=true,
  linkcolor=sncfviolet,
  citecolor=vertvalidation,
  urlcolor=sncfrouge,
  pdftitle={\TitreRapport},
  pdfauthor={\Auteur}
}

\pagestyle{fancy}
\fancyhf{}
\fancyhead[L]{\small Assistant GraphRAG Arianes}
\fancyhead[R]{\small\nouppercase{\leftmark}}
\fancyfoot[C]{\thepage}
\setlength{\headheight}{15pt}

\lstset{
  backgroundcolor=\color{fondcode},
  basicstyle=\ttfamily\small,
  frame=single,
  rulecolor=\color{black!20},
  breaklines=true,
  columns=fullflexible,
  showstringspaces=false,
  numbers=left,
  numberstyle=\tiny\color{sncfgris},
  keywordstyle=\color{sncfviolet}\bfseries,
  commentstyle=\color{vertvalidation},
  stringstyle=\color{sncfrouge}
}

\newcommand{\code}[1]{\texttt{#1}}
\newcommand{\etat}[1]{\textcolor{sncfviolet}{\textbf{#1}}}
\newcommand{\attention}[1]{%
  \begin{center}\fcolorbox{sncfrouge}{sncfrouge!4}{%
  \begin{minipage}{0.9\textwidth}\textbf{Point de vigilance.} #1\end{minipage}}\end{center}}
\newcommand{\resultat}[1]{%
  \begin{center}\fcolorbox{vertvalidation}{vertvalidation!4}{%
  \begin{minipage}{0.9\textwidth}\textbf{Résultat.} #1\end{minipage}}\end{center}}

% Si l'image n'existe pas encore, Overleaf affiche une zone réservée.
\newcommand{\capture}[3]{%
\begin{figure}[H]
  \centering
  \IfFileExists{figures/#1}{%
    \includegraphics[width=0.93\textwidth,height=0.68\textheight,keepaspectratio]{figures/#1}%
  }{%
    \fbox{\begin{minipage}[c][5cm][c]{0.88\textwidth}
      \centering\Large Capture à ajouter\\[0.4em]
      \normalsize\texttt{figures/#1}\\[0.5em]
      \small #2
    \end{minipage}}%
  }
  \caption{#2}
  \label{#3}
\end{figure}}

\tikzset{
  bloc/.style={draw=sncfgris,rounded corners,fill=white,minimum height=1cm,
    text width=3.1cm,align=center,font=\small},
  source/.style={bloc,fill=sncfviolet!8},
  traitement/.style={bloc,fill=sncfrouge!6},
  stockage/.style={bloc,fill=vertvalidation!8},
  fleche/.style={-{Latex[length=2.4mm]},thick,draw=sncfgris}
}
